% CLAUDE REPORT::WCDERIV-0001::2026-11-12T05:51Z
% !TEX program = xelatex
% Compiles with XeLaTeX or pdfLaTeX (ASCII source).
\documentclass[11pt]{article}
\usepackage[margin=1in]{geometry}
\usepackage{amsmath,amssymb,amsthm}
\usepackage{xurl}
\newtheorem{theorem}{Theorem}
\newtheorem{lemma}[theorem]{Lemma}
\newtheorem{proposition}[theorem]{Proposition}
\newtheorem{corollary}[theorem]{Corollary}
\theoremstyle{definition}
\newtheorem{definition}[theorem]{Definition}
\theoremstyle{remark}
\newtheorem*{remark}{Remark}
\newcommand{\SAT}{\mathrm{SAT}}
\newcommand{\UNSAT}{\mathrm{UNSAT}}
\newcommand{\Wb}{\overline W}
\newcommand{\Sound}{\mathrm{Sound}}
\newcommand{\status}[1]{\textbf{[#1]}}

\title{CLAUDE WCDERIV-0001\\[2pt]
\large Worst-case minimum derivation cost: the optimal choice of the fixed system $R$,\\
and an exact sandwich for the weakest bridge B3}
\author{CLAUDE (scheduled routine, shared worst-case directive)}
\date{2026-11-11 21:51 America/Los\_Angeles}

\begin{document}
\maketitle

\noindent\textbf{Status.} P versus NP remains \textbf{OPEN}. This run proves no separation. It
makes the user's ``free optimal reduction sequence'' program exact at its best possible choice
of fixed system, and shows where its two remaining obligations sit. Additive report: no schedules,
role allocations, ownership regions, or root \texttt{main.tex} were changed. This is not an LSPV
run, so it writes no ownership region. Living log:
\url{research/agents/claude/LIVING_LOG.md}.

\section{Setup (as in the five-role directive of 2026-11-02 12:04)}

Fix a Cook--Reckhow refutation system $R$, i.e.\ a polynomial-time verifier $V_R(F,\pi)$ that is
sound and complete for $\UNSAT$. Take $\operatorname{cost}_R(\pi)=|\pi|$, or any weighting
that is polynomially related to it. Define
\[
W_R(F)=\min_{\pi:\,V_R(F,\pi)=1}\operatorname{cost}_R(\pi),\qquad
\Wb_R(N)=\max_{F\in\UNSAT,\ |\langle F\rangle|\le N}W_R(F),
\]
together with the two obligations
\[
\mathrm{LB}(R):\ \Wb_R(N)\ne N^{O(1)},\qquad
\mathrm{B3}(R):\ \SAT\in P\ \Longrightarrow\ \Wb_R(N)\le N^{O(1)}.
\]
$\mathrm{LB}(R)\wedge\mathrm{B3}(R)$ implies $\SAT\notin P$ \status{PROVED/EXACT, trivial}.

\paragraph{Multisign DAG hypergraph reading.} The minimal reduced logical state of an
unsatisfiable 3-CNF hypergraph is the empty signed hyperedge $\bot$. The minimal work needed to
reach it, minimized over all valid reduction DAGs with shared nodes, is $W_R(F)$ for the calculus $R$
whose moves are the declared hypergraph reductions. Granting free knowledge of the best sequence is
legitimate for a \emph{lower} bound. The inner minimum is a nondeterministic certificate length,
so any search for the sequence only adds cost. \status{PROVED/EXACT}. The same freedom is exactly
what makes B3 nontrivial. A lower bound on certificate length bounds an algorithm only if the
algorithm's run can be \emph{converted into} a short certificate, and that conversion is B3.

\section{Monotonicity along p-simulation}

\begin{lemma}[\status{PROVED/EXACT}]\label{lem:mono}
If $S\le_p R$ ($R$ p-simulates $S$), then $\mathrm{B3}(S)\Rightarrow\mathrm{B3}(R)$ and
$\mathrm{LB}(R)\Rightarrow\mathrm{LB}(S)$.
\end{lemma}
\begin{proof}
$W_R(F)\le q(|F|+W_S(F))$ for a fixed polynomial $q$, so $\Wb_R(N)\le q(N+\Wb_S(N))$. Both
claims follow.
\end{proof}

So B3 gets easier, and LB harder, as $R$ moves up the simulation order. The two obligations pull in
opposite directions.

\section{The canonical extremal system $R_T$}

\begin{definition}
Fix a recursively axiomatized theory $T$ (e.g.\ PA or ZFC) that is sound for $\Pi^0_1$ sentences.
An $R_T$-refutation of $F$ is a tuple $\langle S,c,\sigma,\pi,1^{m}\rangle$, where:
\begin{itemize}
\item $S$ is a verifier code;
\item $\sigma$ is a $T$-proof of the $\Pi^0_1$ sentence $\Sound(S,c)$, ``for all $G,\rho$: if $S$
accepts $(G,\rho)$ within $(|G|+|\rho|+2)^c$ steps, then $G$ is unsatisfiable'';
\item $S$ accepts $(F,\pi)$ within $(|F|+|\pi|+2)^c$ steps;
\item $m=(|F|+|\pi|+2)^c$.
\end{itemize}
\end{definition}

The padding $1^m$ makes verification linear-time in the proof length. Checking $\sigma$ is
polynomial in $|\sigma|$. Soundness of $R_T$ follows from the $\Pi^0_1$-soundness of $T$.
Completeness already holds through the Resolution entry with its constant-size $\sigma$. So $R_T$ is a
Cook--Reckhow system \status{PROVED/EXACT}. This is the standard Krajíček--Pudlák construction
\cite{KP89}, restated in the directive's notation.

\begin{theorem}[Extremality \status{KNOWN THEOREM, KP89; restated EXACT}]\label{thm:ext}
For every clocked verifier $(S,c)$ with $T\vdash\Sound(S,c)$, $S\le_p R_T$. The overhead is
the fixed constant $|\langle S,c,\sigma\rangle|$ plus the padding $(|F|+|\pi|+2)^c$. Consequently:
\begin{enumerate}
\item B3 for any $T$-provably sound system implies $\mathrm{B3}(R_T)$;
\item $\mathrm{LB}(R_T)$ implies LB for every $T$-provably sound system. This includes Resolution,
Cutting Planes, Polynomial Calculus, bounded-depth Frege, Frege, Extended Frege, and the
historical mixed calculus $R_{\mathrm{mix341}}$, provided its soundness is $T$-provable.
\end{enumerate}
\end{theorem}

So $R_T$ is the \emph{optimal} fixed system for the user's program among $T$-sound systems. It
minimizes the bridge obligation, and every other admissible choice either has a harder bridge or a
lower bound already implied by $\mathrm{LB}(R_T)$.

\begin{theorem}[Exact B3 sandwich \status{PROVED/EXACT}]\label{thm:sand}
Let $T^{+}=T+\text{``}T\text{ is }\Pi^0_1\text{-sound''}$. Then
\[
\bigl[\SAT\in P\Rightarrow\exists\ T\text{-provably sound p-bounded }S\bigr]
\ \Longrightarrow\ \mathrm{B3}(R_T)\ \Longrightarrow\
\bigl[\SAT\in P\Rightarrow\exists\ T^{+}\text{-provably sound p-bounded }S\bigr].
\]
\end{theorem}
\begin{proof}
The left arrow is Theorem \ref{thm:ext}(1). For the right arrow, $T^{+}$ proves
$\Sound(R_T,c_0)$ for the linear clock $c_0$ of $R_T$, so $S=R_T$ is a witness.
\end{proof}

\begin{corollary}[\status{PROVED/EXACT}]\label{cor:alg}
If some polynomial-time algorithm $A$ (clock $n^k$) decides $\SAT$ and $T\vdash$ ``$A$ is correct'',
then $\Wb_{R_T}(N)\le O(N^{k'})$. Contrapositive:
$\mathrm{LB}(R_T)\Rightarrow$ no polynomial-time SAT algorithm has $T$-provable correctness.
\end{corollary}
\begin{proof}
Let $S_A(F,1^{|F|^k})$ accept iff $A$ outputs UNSAT within $|F|^k$ steps. Then $T$ proves
$\Sound(S_A,k)$, $S_A$ is p-bounded, and Theorem \ref{thm:ext} applies.
\end{proof}

\paragraph{Exact residual obligations at the optimal choice.}
\begin{itemize}
\item \textbf{LB-side.} $\mathrm{LB}(R_T)$ \status{OPEN}. By Theorem \ref{thm:ext}(2), it implies
superpolynomial Extended Frege lower bounds. No such bound is known: only small polynomial
lower bounds are known for Frege and EF, and the 2025--2026 literature surveyed below gives
none. It is also stronger than ``$R_T$ is not p-optimal'' type statements tied to Pudlák's
conjectures \cite{Pud17}.
\item \textbf{Bridge-side.} The \emph{provability gap} \status{OPEN}. This is the world where
$\SAT\in P$ but every p-bounded system, including every correct polynomial SAT decider viewed as a
system, has soundness unprovable in $T$ (and in $T^{+}$). By Theorem \ref{thm:sand}, B3 fails for
$R_T$ only inside this world. By Lemma \ref{lem:mono} and Theorem \ref{thm:ext}, \emph{no other
$T$-provably sound choice of $R$ can do better}. The gap cannot be removed by changing the
calculus. It needs an argument about formal provability of correctness of fast SAT algorithms.
\end{itemize}

This is consistent with, and sharpens, the corpus finding (LSPV-5/376, LSPV-1/337). That finding
falsifies a bridge requiring \emph{every} correct decider to be $T$-provable, but leaves open the
existential version. Theorem \ref{thm:sand} locates the existential version between two exact
provability statements.

\section{Where the two halves cross in the hypergraph hierarchy}

\begin{proposition}[\status{PROVED/EXACT} from known theorems]\label{prop:cross}
Suppose the reductions of the multisign DAG hypergraph are hyperedge resolution, weakening,
subsumption and node sharing, with no new variables. Then $R_{\mathrm{hyp}}\le_p\mathrm{Res}$.
Hence $\mathrm{LB}(R_{\mathrm{hyp}})$ holds unconditionally: $2^{\Omega(n)}$ for random 3-CNF
\cite{CS88} and for Tseitin formulas on expanders \cite{Urq87}, both with $N=\Theta(n\log n)$. As a
result, $\mathrm{B3}(R_{\mathrm{hyp}})\Leftrightarrow\SAT\notin P$.
\end{proposition}

The last equivalence repeats the corpus firewall LSPV-1/398 for Resolution; it is \emph{not}
progress. Adding auxiliary signed hyperedges that define new variables (extension) makes the
calculus equivalent to Extended Resolution, which is p-equivalent to EF. Here LB is OPEN, while
B3 is strictly easier (Lemma \ref{lem:mono}). The two obligations thus trade places between
Resolution and EF. No currently known system has a proved LB together with a proved or plausible
B3: weak systems have $\mathrm{B3}\equiv$ the full problem, and strong systems have LB beyond the
state of the art. \status{PROVED/EXACT as a classification of known systems; no claim about
unknown intermediate systems.}

\section{Test of the directive's priority 4 (alternative reasoning)}

The hard families of Proposition \ref{prop:cross} do \emph{not} survive stronger reasoning, so a
hypergraph-only lower bound must never be reported as an $R_T$ bound:
\begin{itemize}
\item Tseitin formulas have polynomial-size refutations in Polynomial Calculus over
$\mathbb{F}_2$ (Gaussian elimination) and in Frege;
\item PHP has polynomial Cutting Planes refutations;
\item random 3-CNF is OPEN for Frege.
\end{itemize}
\status{KNOWN THEOREM; hostile control passed}: the report does not transfer them.

\section*{Barrier-transfer literature preflight}

The obstacles are (i) superpolynomial lower bounds for EF-or-stronger, and (ii) the provability
gap. Classification:
\begin{itemize}
\item \textbf{Directly transferable theorem:} Krajíček--Pudlák reflection/simulation
\cite{KP89}, and Cook's PV / EF correspondence \cite{Coo75}.
\item \textbf{Useful analogy needing a new theorem:} Gödel-style obstruction theory. The finite
reflection separator of LSPV-1/337 shows the gap is real for \emph{specific} deciders. A transfer
theorem would have to show that \emph{some} correct fast decider is $T$-provable whenever any is.
\item \textbf{Misleading or failed transfer:} the PQ ledger item 5 (black-box extraction of the
resource from an arbitrary solver). An answer-only black box yields no certificate for UNSAT
answers, so it gives no B3 without the provability gap being closed. Classified FAILED TRANSFER.
Items 1--4 and 6--7 are not relevant this run.
\end{itemize}

\section*{Literature checked this run (2025--2026)}
\begin{itemize}
\item Egidy--Glaßer (STOC 2025) build oracles showing that positive answers on the existence of
p-optimal proof systems for SAT/TAUT-type sets cannot be proved by relativizing methods. So
relativizing proofs cannot settle the existence question that Theorem \ref{thm:sand} sits next to.
\item Monroe (arXiv 2507.13576), on characterizing p-simulation between theories. Version~3
\emph{withdraws} the claim that its Theorem~3.3 implies that no p-optimal system exists. Flagged
as a \textbf{stale claim}: do not cite the earlier version.
\item Alekseev, ``A lower bound for Polynomial Calculus with extension rule'' (Theory of Computing
22, 2026). The system simulates EF, but per its own discussion this does not yield an EF lower bound
for CNF tautologies. It is not an $\mathrm{LB}(R_T)$ input.
\item ``The Proof Analysis Problem'' (FOCS 2025; arXiv 2506.16956) shows that analysing EF proofs is
NP-complete. This is a search or verification fact, irrelevant to the free-sequence minimum.
\item Fleming--Grosser--Pitassi--Robere, ``Provable reductions in TFNP'' (ECCC TR26-109).
EF-based systems for TFNP reductions; not an $\mathrm{LB}(R_T)$ input.
\end{itemize}
None of these establishes a superpolynomial EF or Frege lower bound.

\section*{CORPUS REVIEWED}
\begin{itemize}
\item \texttt{AGENTS.md}; the five-role worst-case directive
(\texttt{research/coordination/2026-11-02\_1204\_WORST\_CASE\_MINIMUM\_DERIVATION\_FIVE\_ROLE\_REFOCUS.md});
the recent knot/link handoff (noted, not used).
\item Codex reports 2026-11-11 1203, 1331, 2017 and 2035: the B3 assumption audit, transfer
recovery and rate amplification, and adaptive collision and compatibility.
\item LSPV-1/336, /337, /354, /364 and /398 (clocked optimality, finite reflection, ER-only
route, quantifier firewall, Resolution back-translation firewall), and LSPV-5/376 (strict
reflection handoff).
\item Prior CLAUDE reports (world ledger, TPORT, ONETAPE).
\end{itemize}
Repository-wide searches covered B3 equivalence statements, ZFC and $T$-provability, and
Haken/Tseitin.

\section*{DUPLICATION / CONTRADICTION AUDIT}
\begin{itemize}
\item \emph{Duplicates:} Proposition \ref{prop:cross}'s equivalence for Resolution is
LSPV-1/398. Corollary \ref{cor:alg} is the KP89/Cook reflection idea already used in
LSPV-1/336--361, here applied to the canonical $R_T$.
\item \emph{New in this run:} the extremality framing of $R_T$ as the optimal fixed $R$ for the
directive (Lemma \ref{lem:mono} + Theorem \ref{thm:ext}) and the two-sided sandwich
(Theorem \ref{thm:sand}). Repository search found no prior statement of the $T$/$T^{+}$ sandwich.
\item No contradictions were found with the corpus.
\end{itemize}

\section*{WHAT CHANGED RELATIVE TO ROOT MAIN}
Nothing in root \texttt{main.tex} was edited; this is a non-LSPV, additive report. Suggested
inputs for LSPV-5 synthesis: the sandwich (Theorem \ref{thm:sand}) and the crossing
classification (Proposition \ref{prop:cross}).

\section*{ADOPTED / REJECTED / RETIRED PRIOR COMPONENTS}
\begin{itemize}
\item \emph{Adopted:} the directive's definitions; LSPV-1/398; and LSPV-5/376's scope
(the existential B3 is not falsified).
\item \emph{Rejected:} treating any hypergraph-only (Resolution-strength) lower bound as
evidence for B3 progress.
\item \emph{Retired:} none.
\end{itemize}

\section*{STALE OR OVERSTATED CLAIMS FOUND}
Monroe arXiv 2507.13576 versions~1--2 (the claim of no p-optimal system, withdrawn in v3). No
overstated claim was found in the corpus files reviewed.

\section*{Status table}
\begin{tabular}{ll}
Free-sequence lower-bound principle & PROVED/EXACT\\
Lemma \ref{lem:mono} (monotonicity) & PROVED/EXACT\\
Theorem \ref{thm:ext} (extremality of $R_T$) & KNOWN THEOREM (KP89), restated EXACT\\
Theorem \ref{thm:sand} (B3 sandwich) & PROVED/EXACT\\
Corollary \ref{cor:alg} & PROVED/EXACT\\
Proposition \ref{prop:cross} & PROVED/EXACT from known theorems\\
$\mathrm{LB}(R_T)$ & OPEN\\
$\mathrm{B3}(R_T)$ (provability gap) & OPEN\\
P versus NP & OPEN
\end{tabular}

\section*{Single next target}
Decide whether the provability gap can be shrunk \emph{existentially}. Given a correct
polynomial-time SAT decider $A$ with unknown correctness proof, construct, uniformly from $A$'s
code, a polynomial-time decider $A'$ and a $T^{+}$-proof of its correctness, or find an oracle or
model obstruction to doing so. Egidy--Glaßer-style oracles are the first hostile control. A
positive result would make B3 for $R_T$ unconditional relative to $T^{+}$. That would leave the
user's program exactly equivalent to $\mathrm{LB}(R_T)$: a superpolynomial lower bound for the
strongest provably sound proof system.

\begin{thebibliography}{9}
\bibitem{KP89} J.~Krajíček, P.~Pudlák. Propositional proof systems, the consistency of first
order theories and the complexity of computations. \emph{J.\ Symb.\ Logic} 54(3), 1989.
\bibitem{Coo75} S.~Cook. Feasibly constructive proofs and the propositional calculus. STOC 1975.
\bibitem{Pud17} P.~Pudlák. Incompleteness in the finite domain. \emph{Bull.\ Symb.\ Logic} 23(4), 2017.
\bibitem{CS88} V.~Chvátal, E.~Szemerédi. Many hard examples for resolution. \emph{J.\ ACM} 35, 1988.
\bibitem{Urq87} A.~Urquhart. Hard examples for resolution. \emph{J.\ ACM} 34, 1987.
\bibitem{EG25} F.~Egidy, C.~Glaßer. Oracle constructions concerning p-optimal proof systems. STOC 2025 (title abbreviated; verify before citing).
\end{thebibliography}

\end{document}
% END_OF_VERIFIED_REPORT::CLAUDE::WCDERIV-0001
